\section{Implementación}
\label{sec:implementacion}

\subsection{Descomposición del trabajo}

El trabajo de una configuración $c$ son $N$ escenarios independientes. El escenario local $j\in[0,N)$ se identifica con el índice global $i = c\cdot N + j$. Como el resultado de cada escenario depende solo de $i$ (Sección~\ref{sec:metodologia}), repartir trabajo equivale a repartir rangos de índices, y cualquier reparto produce el mismo resultado. Cada backend acumula por configuración tres cantidades enteras: $k$ (inundaciones), $n$ (escenarios simulados) y el histograma de 512 clases. La reducción, entre hilos o entre procesos, es una suma de enteros sin signo de 64 bits; con \codigo{uint64\_t} se evita el desbordamiento que un \codigo{int} sufriría a partir de $2{,}1\times10^{9}$ escenarios.

Todos los ejecutables comparten la misma línea de comandos, el mismo archivo de configuración y el mismo formato CSV de resultados, que incluye $\phat$, el intervalo de Wilson, el histograma completo, los metadatos de hardware y el commit.

\subsection{Versión secuencial}

\codigo{pp\_serial} recorre $j = 0,\dots,N-1$, llama a \codigo{simular\_escenario} y acumula. Es la referencia de V1 y la línea base $T_1$ de todas las métricas.

\subsection{OpenMP}

\codigo{pp\_openmp} procesa el intervalo $[j_0, j_1)$ con un único bucle paralelo:
\begin{center}
\small
\codigo{\#pragma omp parallel for schedule(runtime) \textbackslash}\\
\codigo{\ \ reduction(+:k,n) reduction(+:hist[:512])}
\end{center}
Cada hilo tiene copias privadas de $k$, $n$ y de un histograma de \qty{4}{\kilo\byte}, y el runtime las suma al cerrar la región; el cuerpo del bucle no escribe en memoria compartida, por lo que no hay carreras (V5). \codigo{schedule(runtime)} permite comparar \codigo{static}, \codigo{dynamic} y \codigo{guided} sin recompilar: el tipo y el tamaño de bloque se fijan con \codigo{omp\_set\_schedule} y el valor efectivo se lee con \codigo{omp\_get\_schedule} y se registra. La reducción se hace por bloque de una sola configuración; reducir los acumuladores de todas las configuraciones a la vez crearía una copia privada de decenas de megabytes por hilo.

\subsection{MPI: dos esquemas algorítmicos}

Los dos ejecutables MPI comparten el flujo: \codigo{MPI\_Init\_thread} con \codigo{MPI\_THREAD\_FUNNELED}; \codigo{MPI\_Barrier} y \codigo{MPI\_Wtime} para iniciar la medición; ejecución del esquema; una única \codigo{MPI\_Reduce} con \codigo{MPI\_SUM} sobre un vector de $2 + 512$ valores \codigo{MPI\_UINT64\_T}; un \codigo{MPI\_Gather} con el resumen de cada rango para comprobar la cobertura (V3); y \codigo{MPI\_Bcast} del código de salida.

\textbf{E1, SPMD estático con reducción (\emph{map-reduce}).} Cada rango $r$ de $P$ procesa un intervalo contiguo de tamaño $\lfloor N/P\rfloor$ o $\lfloor N/P\rfloor + 1$: los primeros $N \bmod P$ rangos reciben un escenario más, sin solapes ni huecos. Durante el cómputo no hay mensajes punto a punto; la comunicación se reduce a una colectiva. Su limitación es que, con nodos heterogéneos, el más lento fija el tiempo total.

\textbf{E2, maestro--trabajador dinámico.} Con $P > 1$, el rango 0 actúa solo como coordinador y mantiene una cola de lotes de $B = \num{62500}$ escenarios. Cada trabajador envía una solicitud (\codigo{REQ}), recibe un rango $[i_0, i_1)$ (\codigo{WORK}) y repite hasta recibir \codigo{STOP}. El coordinador envía exactamente un \codigo{STOP} por trabajador y lleva la cuenta de trabajadores activos, lo que excluye interbloqueos; al terminar verifica que los lotes asignados cubren $[0, N)$ sin solapes. El esquema equilibra la carga de forma automática ---un nodo más rápido pide más lotes---, al costo de un mensaje por lote y de dedicar el rango 0 a coordinar, en línea con los marcos de planificación dinámica para clústeres heterogéneos~\cite{tse2010}.

\subsection{Híbrido MPI + OpenMP}

\codigo{pp\_hibrido} reparte $[0, N)$ entre procesos como E1 y, dentro de cada proceso, ejecuta el bucle OpenMP con $T$ hilos sobre su intervalo; el presupuesto de recursos es $P\times T$. Todas las llamadas MPI ocurren antes o después de la región paralela y solo desde el hilo principal, por lo que basta el nivel \codigo{MPI\_THREAD\_FUNNELED}; si el nivel concedido es menor, el programa termina con error. Open MPI ancla por defecto cada proceso a un núcleo, de modo que los $T$ hilos compartirían ese núcleo; por eso el híbrido se lanza con \codigo{-{}-bind-to none} y \codigo{OMP\_PROC\_BIND=close}, \codigo{OMP\_PLACES=cores}, o bien con \codigo{-{}-map-by ppr:1:node:PE=$T$}, que reserva $T$ núcleos por proceso.
% TODO-T33: citar aquí el efecto medido de --bind-to none con el CSV versionado de T27 (results/bench/hibrido.csv); la cifra exploratoria previa no está versionada y no se publica.

\subsection{CUDA: decisión NO-GO}

Monte Carlo es un caso típico de cómputo en GPU: un hilo por escenario, aritmética regular y una reducción final~\cite{benguigui2012}. El diseño la preveía como extensión desacoplada (EXT1: kernel con reducción en memoria compartida, V2, V8 y barrido de bloques X5; EXT2: la GPU como trabajador del esquema E2), condicionada a que el núcleo CPU estuviese completo y a disponer de un entorno CUDA operativo al cierre de la semana~10. Ninguna de las dos condiciones se cumplió: E2, E4 y E5 no tenían aún evidencia completa, y aunque había hardware (RTX~4060 y RTX~5090 Laptop), ningún equipo tenía el CUDA Toolkit instalado, ambas GPU trabajan en modo WDDM ---con el \emph{watchdog} TDR de Windows, de unos 2~s, sobre los kernels largos--- y la imagen \codigo{pp-g8} no incluye \codigo{nvcc}. El equipo decidió no activar CUDA y dedicar ese tiempo a los experimentos CPU. No se publican cifras de GPU, y V2, V8 y X5 quedan fuera de alcance. Como Philox ya es una función pura, marcarla \codigo{\_\_host\_\_ \_\_device\_\_} bastaría para reutilizar el mismo flujo aleatorio en una futura versión GPU.

\subsection{Contenedores, Compose y Swarm}

La imagen \codigo{pp-g8} se construye una vez y es idéntica en todos los nodos. En Compose, un archivo define un coordinador y dos trabajadores en una red común; los binarios se compilan una vez en un volumen compartido, y el coordinador lanza \codigo{mpirun} por SSH con clave entre contenedores. En Swarm, el \emph{stack} define un servicio por equipo físico, restringido por la etiqueta de nodo \codigo{pp\_nodo=1|2|3}, conectado a una red \emph{overlay} con resolución DNS por tarea (\codigo{endpoint\_mode: dnsrr}); el \emph{hostfile} de \codigo{mpirun} se genera a partir de los servicios. El mismo \emph{stack} puede simularse en un único anfitrión para pruebas, sabiendo que esa simulación no sirve para medir escalabilidad. La red universitaria suele bloquear los puertos de la \emph{overlay} (4789/udp, 7946 y 2377), por lo que el despliegue se planificó para una red doméstica o compartida.
% TODO-T33: registrar la ejecución de principio a fin en los tres equipos (logs de T25/T34) y la latencia/ancho de banda de la red usada.
